%------------------------------------------------------------------------------%
\documentclass[fleqn,utf8,aspectratio=169,12pt]{beamer}

\usepackage{gaal}

\title{Ângulo entre Vetores}

%------------------------------------------------------------------------------%
\begin{document}

\addtitle

%------------------------------------------------------------------------------%
\section{Ângulo entre Vetores}

%------------------------------------------------------------------------------%
\begin{frame}{Calculando o Ângulo}
  \onslide<+->{Com as componentes dos vetores podemos calcular o ângulo entre eles}

  \onslide<+->{Se nenhum vetor é nulo}
  \begin{align*}
    \onslide<+->{\vec{u}\cdot\vec{v} &  = \norm{\vec{u}} \norm{\vec{v}} \cos(\theta) \\}
    \onslide<+->{\cos\theta & = \frac{\vec{u}\cdot\vec{v}}{\norm{\vec{u}} \norm{\vec{v}}} \\}
    \onslide<+->{\theta & = \arccos \left(
                  \frac{\vec{u}\cdot\vec{v}}{\norm{\vec{u}} \norm{\vec{v}}}
                 \right)}
  \end{align*}
\end{frame}

%------------------------------------------------------------------------------%
\begin{frame}{Dimensões Maiores}
  A  expressão
  \(\;
    \vec{u}\cdot\vec{v}
    \;=\; \norm{\vec{u}} \norm{\vec{v}} \cos(\theta)
  \;\)
  vale em $\R^n$

  \pause
  \bigskip
  Definimos o ângulo entre vetores em dimensões maiores por
  \pause
  \[
    \theta = \arccos \left(
              \frac{\vec{u}\cdot\vec{v}}{\norm{\vec{u}} \norm{\vec{v}}}
             \right)
  \]
  \pause
  desde que
  \(
    \vec{u} \neq \vec{0}
  \)
  e
  \(
    \vec{v} \neq \vec{0}
  \)
\end{frame}

%------------------------------------------------------------------------------%
%------------------------------------------------------------------------------%
\begin{question}
  Determine o ângulo entre os vetores
  \[
    \vec{w} = (1, 1, 0)
    \qquad
    \vec{v} = (1, 0, 1)
  \]
\end{question}

%------------------------------------------------------------------------------%
\begin{resolution}{Produto Escalar}
  \onslide<+->{Produto Escalar}
  \begin{align*}
    \onslide<+->{\vec{w}\cdot\vec{v}}
    \onslide<+->{& = (1, 1, 0) \cdot  (1, 0, 1) \\}
    \onslide<+->{& = 1 \times 1 + 1 \times 0 + 0 \times 1 \\}
    \onslide<+->{& = 1}
  \end{align*}
\end{resolution}

%------------------------------------------------------------------------------%
\begin{resolution}{Norma}
\begin{columns}[t]
\column{0.4\linewidth}
  \onslide<+->{Norma de $\vec{w}$}
  \begin{align*}
    \onslide<+->{\norm{u}}
    \onslide<+->{& = \sqrt{u_1^2 + u_2^2 + u_3^2} \\}
    \onslide<+->{& = \sqrt{1^2 + 1^2 + 0^2} \\}
    \onslide<+->{& = \sqrt{2}}
  \end{align*}
\column{0.5\linewidth}
  \onslide<+->{Norma de $\vec{w}$}
  \begin{align*}
    \onslide<+->{\norm{v}}
    \onslide<+->{& = \sqrt{v_1^2 + v_2^2 + v_3^2} \\}
    \onslide<+->{& = \sqrt{1^2 + 0^2 + 1^2} \\}
    \onslide<+->{& = \sqrt{2}}
  \end{align*}
\end{columns}
\end{resolution}

%------------------------------------------------------------------------------%
\begin{resolution}{Solução}
\begin{columns}[t]
\column{0.5\linewidth}
  \begin{align*}
    \onslide<+->{\cos(\theta)}
    \onslide<+->{& = \frac{u \cdot v}{\norm{u} \, \norm{v}} \\}
    \onslide<+->{& = \frac{1}{\sqrt{2}\sqrt{2}} \\}
    \onslide<+->{& = \frac{1}{2} }
  \end{align*}
\column{0.5\linewidth}
  \onslide<+->{\[
    \theta = \frac{\pi}{3}
  \]}
\end{columns}
\end{resolution}

%------------------------------------------------------------------------------%
\endinput

%------------------------------------------------------------------------------%
\begin{question}
  Determine o ângulo entre os vetores de $\R^5$
  \[
    \vec{u} = (1, 2, -1, 3, 0)
  \]

  \[
    \vec{v} = (2, -1, 4, 0, 1)
  \]
\end{question}

%------------------------------------------------------------------------------%
\begin{resolution}{Produto Escalar}
  \onslide<+->{Produto escalar}
  \begin{align*}
    \onslide<+->{\vec{u} \cdot \vec{v}}
    \onslide<+->{& = (1, 2, -1, 3, 0) \cdot (2, -1, 4, 0, 1) \\}
    \onslide<+->{& = 1 \times 2 + 2(-1) + (-1)4 + 3 \times 0 + 0 \times 1 \\}
    \onslide<+->{& = 2 - 2 - 4 + 0 + 0 \\}
    \onslide<+->{& = -4}
  \end{align*}
\end{resolution}

%------------------------------------------------------------------------------%
\begin{resolution}{Normas}
\begin{columns}
\column{0.5\textwidth}
    \onslide<+->{Norma de $\vec{u}$}
    \begin{align*}
      \onslide<+->{\norm{\vec{u}}}
      \onslide<+->{& = \sqrt{1^2 + 2^2 + (-1)^2 + 3^2 + 0^2} \\}
      \onslide<+->{& = \sqrt{1 + 4 + 1 + 9} \\}
      \onslide<+->{& = \sqrt{15}}
    \end{align*}
\column{0.5\textwidth}
  \onslide<+->{Norma de $\vec{v}$}
  \begin{align*}
    \onslide<+->{\norm{\vec{v}}}
      \onslide<+->{& = \sqrt{2^2 + (-1)^2 + 4^2 + 0^2 + 1^2} \\}
      \onslide<+->{& = \sqrt{4 + 1 + 16 + 1} \\}
      \onslide<+->{& = \sqrt{22}}
  \end{align*}
\end{columns}
\end{resolution}

%------------------------------------------------------------------------------%
\begin{resolution}{Solução}
  \[
    \cos\theta
    \pause
    \;=\; \frac{\vec{u} \cdot \vec{v}}{\norm{\vec{u}}\, \norm{\vec{v}}}
    \pause
    \;=\; \frac{-4}{\sqrt{15}\sqrt{22}}
    \pause
    \;=\; -\frac{4}{\sqrt{330}}
  \]

  \pause
  \bigskip
  \[
    \theta
    \pause
    \;=\; \arccos \left(
        -\frac{4}{\sqrt{330}}
      \right)
    \pause
    {\color{gray}
    \;\approx\; 1,79
    \pause
    \;\approx \ang{60.85}
    }
  \]
  \end{resolution}

%------------------------------------------------------------------------------%
\endinput




%------------------------------------------------------------------------------%
\section{Ortogonalidade}

%------------------------------------------------------------------------------%
\begin{frame}{Ortogonalidade}
  Dois vetores não nulos são \emph{ortogonais}

  \pause
  quando o ângulo entre eles é reto,
  \(
    \theta = \frac{\pi}{2}
  \)
  \pause
  \[
    \cos(\theta)
    \pause
    \;=\; \cos\left(\frac{\pi}{2}\right)
    \pause
    \;=\; 0
  \]
  \[
    \pause
    \vec{u}\cdot\vec{v}
    \pause
    \;=\; \norm{\vec{u}} \norm{\vec{v}} \cos(\theta)
    \pause
    \;=\; \norm{\vec{u}} \norm{\vec{v}} \cos\left(\frac{\pi}{2}\right)
    \pause
    \;=\; 0
  \]

  \pause
  Portanto
  \[
    \vec{u}\perp\vec{v}
    \quad\Leftrightarrow\quad
    \vec{u}\cdot\vec{v} = 0
  \]
\end{frame}

%------------------------------------------------------------------------------%
%------------------------------------------------------------------------------%
\begin{question}
  Determine $k$ para que os vetores
  \[
    \vec{w} = \vetortri{1}{k}{2}
    \qquad
    \vec{v} = \vetortri{3}{2}{-1}
  \]
  sejam ortogonais

  \pause
  \bigskip
  \bigskip
  Condição de ortogonalidade
  \(
    \vec{w} \cdot \vec{v} = 0
  \)
\end{question}

%------------------------------------------------------------------------------%
\begin{resolution}{Solução}
  \begin{align*}
    \onslide<+->{\vec{w} \cdot \vec{v} & = 0 \\}
    \onslide<+->{\vetortri{1}{k}{2} \cdot \vetortri{3}{2}{-1} & = 0 \\}
    \onslide<+->{1 \times 3 + k \times 2 + 2(-1) & = 0 \\}
    \onslide<+->{3 + 2k - 2 & = 0 \\}
    \onslide<+->{2k + 1 & = 0 \\}
    \onslide<+->{k & = -\frac{1}{2}}
  \end{align*}
\end{resolution}

%------------------------------------------------------------------------------%
\endinput


%------------------------------------------------------------------------------%
%\section{Lista Mínima}
%
%%------------------------------------------------------------------------------%
%\begin{frame}{Lista Mínima}
%  \begin{enumerate}
%    \item
%  \end{enumerate}
%
%  \vfill
%  Atenção: A prova é baseada no livro, não nas apresentações
%\end{frame}

%------------------------------------------------------------------------------%
\end{document}
%------------------------------------------------------------------------------%
